\chapter{SYSTEM DESIGN}

\justify

This chapter describes the architecture and design of the HPDC Defect Detection System, covering the full RF-DETR model architecture, data pipeline design, and dataset structure.

\section{Architecture of the System}
\justify

The system architecture consists of three major components: (1) the Data Engineering Pipeline, (2) the RF-DETR Detection Model, and (3) the Evaluation and Inference Engine. RF-DETR Medium uses a DINOv2 ViT-B backbone, a multi-scale projector neck, and a multi-scale deformable DETR decoder.

\begin{figure}[H]
\centering
\begin{tikzpicture}[node distance=0.5cm,
mainblock/.style={rectangle, rounded corners=5pt, draw=magnaDark,
fill=magnaLightGray, text width=9cm,
minimum height=1.1cm, align=center, font=\small},
subblock/.style={rectangle, rounded corners=3pt, draw=magnaGray,
fill=white, text width=8cm,
minimum height=0.9cm, align=center, font=\footnotesize},
arrow/.style={-Stealth, thick, color=magnaRed}]

\node[mainblock, fill=magnaBlue!15, draw=magnaBlue] (input)
    {\textbf{Input Patch} --- 1024$\times$1024 pixels};
\node[mainblock, fill=magnaRed!10, draw=magnaRed, below=0.6cm of input] (backbone)
    {\textbf{BACKBONE: DINOv2-style Vision Transformer (ViT-B)}};
\node[subblock, below=0.2cm of backbone] (pembed)
    {Patch Embedding (16$\times$16) + Position Embeddings + Register Tokens};
\node[subblock, below=0.2cm of pembed] (vitblocks)
    {Transformer Encoder: Windowed Attention (local) $\leftrightarrow$ Global Attention\\
    Alternating pattern --- Output: Multi-scale feature maps \{C3, C4, C5\}};
\node[mainblock, fill=successGreen!10, draw=successGreen, below=0.6cm of vitblocks] (neck)
    {\textbf{MULTI-SCALE PROJECTOR (Neck)}\\C2f-style feature aggregation across resolution levels};
\node[mainblock, fill=warningOrange!10, draw=warningOrange, below=0.6cm of neck] (decoder)
    {\textbf{DEFORMABLE DETR DECODER}};
\node[subblock, below=0.2cm of decoder] (queries)
    {300 Learnable Object Queries};
\node[subblock, below=0.2cm of queries] (deformattn)
    {Multi-scale Deformable Cross-Attention ($K$ sparse sampling points per query)\\
    Feed-Forward Network + Self-Attention between queries};
\node[subblock, below=0.2cm of deformattn] (heads)
    {Classification MLP $\to$ class label \quad|\quad Bounding Box MLP $\to$ $(c_x, c_y, w, h)$};
\node[mainblock, fill=magnaBlue!10, draw=magnaBlue, below=0.6cm of heads] (matching)
    {\textbf{BIPARTITE MATCHING (Training)} --- Hungarian Algorithm\\
    Loss = Classification + $\mathcal{L}_1$ bbox + GIoU};
\node[mainblock, fill=successGreen!15, draw=successGreen, below=0.6cm of matching] (output)
    {\textbf{Output: \{class, confidence, bbox\}} --- No NMS required};

\draw[arrow] (input)      -- (backbone);
\draw[arrow] (backbone)   -- (pembed);
\draw[arrow] (pembed)     -- (vitblocks);
\draw[arrow] (vitblocks)  -- (neck);
\draw[arrow] (neck)       -- (decoder);
\draw[arrow] (decoder)    -- (queries);
\draw[arrow] (queries)    -- (deformattn);
\draw[arrow] (deformattn) -- (heads);
\draw[arrow] (heads)      -- (matching);
\draw[arrow] (matching)   -- (output);
\end{tikzpicture}
\caption{RF-DETR Medium complete architecture. Each patch is processed end-to-end from DINOv2 backbone through deformable decoder to produce detections without NMS.}
\label{fig:rfdetr_arch}
\end{figure}

\section{Class Diagram (with Brief Explanation)}
\justify

The system is organised into four primary Python classes. The \texttt{PatchGenerator} handles sliding-window extraction and Shapely polygon clipping. The \texttt{DatasetBuilder} applies balancing and augmentation. The \texttt{RFDETRTrainer} wraps the RF-DETR training loop with EMA tracking. The \texttt{DefectDetector} handles inference on full images.

\begin{figure}[H]
\centering
\begin{tikzpicture}[
    classbox/.style={rectangle, draw=magnaDark, fill=magnaLightGray,
    minimum width=4.5cm, font=\small},
    arrow/.style={-Stealth, thick, color=magnaBlue}
]
% PatchGenerator
\node[classbox, minimum height=2.2cm, align=left] (pg) at (0,0) {
    \textbf{PatchGenerator}\\
    \rule{4.3cm}{0.4pt}\\
    -- patch\_size: int\\
    -- overlap: float\\
    \rule{4.3cm}{0.4pt}\\
    + generate(image, annotations)\\
    + clip\_polygons(patch\_box)
};
% DatasetBuilder
\node[classbox, minimum height=2.2cm, align=left] (db) at (5.5,0) {
    \textbf{DatasetBuilder}\\
    \rule{4.3cm}{0.4pt}\\
    -- balance\_ratio: float\\
    -- aug\_classes: list\\
    \rule{4.3cm}{0.4pt}\\
    + balance(patches)\\
    + augment\_crack\_patches()
};
% RFDETRTrainer
\node[classbox, minimum height=2.2cm, align=left] (tr) at (0,-3.5) {
    \textbf{RFDETRTrainer}\\
    \rule{4.3cm}{0.4pt}\\
    -- model: RFDETRMedium\\
    -- epochs: int\\
    \rule{4.3cm}{0.4pt}\\
    + train(dataset)\\
    + save\_ema\_checkpoint()
};
% DefectDetector
\node[classbox, minimum height=2.2cm, align=left] (dd) at (5.5,-3.5) {
    \textbf{DefectDetector}\\
    \rule{4.3cm}{0.4pt}\\
    -- checkpoint\_path: str\\
    -- conf\_threshold: float\\
    \rule{4.3cm}{0.4pt}\\
    + detect(image\_path)\\
    + annotate\_image(detections)
};

\draw[arrow] (pg) -- (db);
\draw[arrow] (db) -- (tr);
\draw[arrow] (tr) -- (dd);
\draw[arrow] (pg) -- (tr);
\end{tikzpicture}
\caption{Simplified class diagram of the four main system components.}
\label{fig:class_diagram}
\end{figure}

\section{Sequence Diagram (with Brief Explanation)}
\justify

The sequence diagram below illustrates the inference flow: a quality engineer submits an image, the system generates patches, runs RF-DETR on each patch, and returns annotated detections.

\begin{figure}[H]
\centering
\begin{tikzpicture}[font=\small]
% Lifelines
\draw[thick] (0,0) -- (0,-7); \node at (0,0.3) {\textbf{Engineer}};
\draw[thick] (3,0) -- (3,-7); \node at (3,0.3) {\textbf{DefectDetector}};
\draw[thick] (6,0) -- (6,-7); \node at (6,0.3) {\textbf{PatchGenerator}};
\draw[thick] (9.5,0) -- (9.5,-7); \node at (9.5,0.3) {\textbf{RF-DETR}};

% Activations
\fill[magnaBlue!20] (2.85,-0.5) rectangle (3.15,-6.5);
\fill[magnaBlue!20] (5.85,-1.2) rectangle (6.15,-2.8);
\fill[magnaBlue!20] (9.35,-3.0) rectangle (9.65,-5.0);

% Messages
\draw[-Stealth, magnaRed] (0,-0.8) -- (3,-0.8) node[midway, above, font=\tiny] {detect(image\_path)};
\draw[-Stealth, magnaRed] (3,-1.4) -- (6,-1.4) node[midway, above, font=\tiny] {generate\_patches(image)};
\draw[-Stealth, dashed, magnaBlue] (6,-2.6) -- (3,-2.6) node[midway, above, font=\tiny] {return patches[ ]};
\draw[-Stealth, magnaRed] (3,-3.2) -- (9.5,-3.2) node[midway, above, font=\tiny] {forward(patch)};
\draw[-Stealth, dashed, magnaBlue] (9.5,-4.8) -- (3,-4.8) node[midway, above, font=\tiny] {return \{class, conf, bbox\}};
\draw[-Stealth, dashed, magnaBlue] (3,-6.2) -- (0,-6.2) node[midway, above, font=\tiny] {return annotated\_image};
\end{tikzpicture}
\caption{Sequence diagram for the defect detection inference flow.}
\label{fig:sequence_diagram}
\end{figure}

\section{ER Diagram and Schema (if applicable)}
\justify

The system does not use a relational database. Data is stored as COCO-format JSON files and PNG image patches on the filesystem. The logical schema is as follows:

\begin{table}[H]
\centering
\caption{COCO JSON Logical Schema}
\begin{tabular}{lll}
\toprule
\rowcolor{tableHeader}\tableheadfont Entity & \tableheadfont Key Fields & \tableheadfont Relationship \\
\midrule
Image           & id, file\_name, width, height & 1:N with Annotations \\
\rowcolor{tableRowAlt}Annotation      & id, image\_id, category\_id, bbox, segmentation & N:1 with Image \\
Category        & id, name (Crack / Cold Shut) & 1:N with Annotations \\
\bottomrule
\end{tabular}
\end{table}

\section{State Transition Diagram (if applicable)}
\justify

The training process transitions through well-defined states as shown below:

\begin{figure}[H]
\centering
\begin{tikzpicture}[
    state/.style={draw, rounded corners, fill=magnaLightGray, minimum width=3cm, minimum height=0.8cm, font=\small, align=center},
    arrow/.style={-Stealth, thick, color=magnaRed}
]
\node[state, fill=magnaBlue!20] (init) at (0,0) {Initialise\\Model};
\node[state] (load) at (4,0) {Load\\Dataset};
\node[state] (train) at (8,0) {Training\\Epoch};
\node[state] (eval) at (8,-2) {Validate\\(mAP)};
\node[state, fill=successGreen!20] (save) at (4,-2) {Save EMA\\Checkpoint};
\node[state, fill=magnaRed!20] (done) at (0,-2) {Training\\Complete};

\draw[arrow] (init)  -- (load);
\draw[arrow] (load)  -- (train);
\draw[arrow] (train) -- (eval);
\draw[arrow] (eval)  -- node[right, font=\tiny] {mAP improved} (save);
\draw[arrow] (eval)  -- node[above, font=\tiny, text width=2cm, align=center] {epoch $<$ max} (train);
\draw[arrow] (save)  -- (done);
\draw[arrow] (save)  |- node[below, font=\tiny] {continue training} (train);
\end{tikzpicture}
\caption{State transition diagram for the RF-DETR training process.}
\end{figure}

\section{Dataset Description}
\justify

\begin{table}[H]
\centering
\caption{Original Dataset Summary}
\begin{tabular}{lc}
\toprule
\rowcolor{tableHeader}\tableheadfont Parameter & \tableheadfont Value \\
\midrule
Dataset name       & Cast\_Cracks\_Annotated\_Coco \\
\rowcolor{tableRowAlt}Training images    & 173 \\
Validation images  & 50 \\
\rowcolor{tableRowAlt}Test images        & 33 \\
\textbf{Total images} & \textbf{256} \\
\rowcolor{tableRowAlt}Image resolution   & 4096 $\times$ 3000 pixels \\
Annotation format  & COCO JSON (polygon-level) \\
\rowcolor{tableRowAlt}Defect classes     & Crack, Cold Shut \\
\bottomrule
\end{tabular}
\end{table}

\begin{table}[H]
\centering
\caption{Patch Dataset Summary (Final --- Experiment 7)}
\begin{tabular}{lcc}
\toprule
\rowcolor{tableHeader}\tableheadfont Split & \tableheadfont Images & \tableheadfont Annotations \\
\midrule
Train       & 2,939 & 2,743 \\
\rowcolor{tableRowAlt}Validation  & 750   & 705 \\
Test        & 495   & 596 \\
\rowcolor{tableRowAlt}\textbf{Total} & \textbf{4,184} & \textbf{4,044} \\
\bottomrule
\end{tabular}
\end{table}

Key EDA findings: the median crack bounding box occupies $\approx$0.065\% of the full image area (roughly $90\times90$ pixels on a 12.3~MP image), confirming that patch-based training is essential. Cold Shut objects are approximately $2\times$ larger than cracks on average, explaining the consistent per-class performance gap.
